\documentclass[11pt,a4paper]{article}

\usepackage{fontspec}
\usepackage{amsmath,amssymb,amsthm}
\usepackage[margin=1in]{geometry}
\usepackage{microtype}
\usepackage{fancyvrb}
\usepackage[hidelinks]{hyperref}
\usepackage{enumitem}

% Lean source uses ℚ, ∑, ∈, →, so the mono font must cover them.
\IfFontExistsTF{DejaVu Sans Mono}
  {\setmonofont{DejaVu Sans Mono}[Scale=0.85]}
  {\IfFontExistsTF{Menlo}{\setmonofont{Menlo}[Scale=0.85]}{}}

\newtheorem{theorem}{Theorem}
\theoremstyle{definition}
\newtheorem{definition}[theorem]{Definition}

\newcommand{\lean}[1]{\texttt{#1}}

\title{Physarum-Inspired Branch-Flow Search\\
\large The machine-checked core of a self-organizing game-tree search}
\author{}
\date{}

\begin{document}
\maketitle

\begin{abstract}
\noindent
Branch-flow search assigns computational traffic to complete root-to-frontier
branches of a game tree, and lets useful traffic reinforce every edge on the
branch it travelled. The substrate is a conductivity network borrowed from the
slime mold \emph{Physarum polycephalum}. A move policy supplies only the
initial conductivities, and no learned action decides where the next unit of
computation goes.

This document states what has been proved about that algorithm and nothing
else. The Lean 4 development contains nineteen theorems. Fifteen are checked by
the Lean kernel against its three standard axioms. Four are exact-rational
computations discharged by \lean{native\_decide}, which also trusts the Lean
compiler, and they are marked wherever they appear. Section~\ref{sec:notproved}
lists what the formalization does not answer, including every question about
playing strength.
\end{abstract}

\section{Scope}

The repository holds one search algorithm. A separate line of work on learned
computation-budget allocation is not included here.

The theorems cover the algebra of the search dynamics: how conductivities are
initialized from a policy, how branch traffic is credited, how unused routes
decay, and how cost accumulates. They hold for any flow assignment and any
usefulness signal meeting the stated sign conditions. That generality is what
makes them cheap to state, and it is also what limits them. No theorem here
says the search plays well.

\section{Where the idea comes from}

\emph{Physarum polycephalum} is an acellular slime mold. It feeds as one
multinucleate cell spread into a network of protoplasmic tubes, and it has no
brain and no central controller. It still solves shortest-path problems, by a
local rule: tubes carrying a large flux of protoplasm thicken, and tubes
carrying little flux thin out and disappear.

Nakagaki, Yamada and T\'oth showed the organism connecting two food sources
through the shortest route in a maze~\cite{nakagaki2000}. Tero, Kobayashi and
Nakagaki then modelled the mechanism as a conductivity network, where each tube
has a conductivity $D_e$ obeying
\begin{equation}
  \frac{dD_e}{dt} = f(|Q_e|) - D_e ,
\end{equation}
with $Q_e$ the flux through the tube under a Kirchhoff flow and $f$
increasing~\cite{tero2007}. Later work reproduced the Tokyo rail network from
the same rule~\cite{tero2010}.

What is borrowed here is the feedback loop,
\begin{equation}
  \text{conductivity} \rightarrow \text{flow} \rightarrow \text{reinforcement}
  \rightarrow \text{conductivity} .
\end{equation}
The network is never designed. It is what the traffic left behind.

\section{What changes when the substrate is a game tree}

A slime mold reinforces a tube in proportion to the flux through it. Flux alone
is the wrong signal for game-tree search, because a search route can carry
heavy traffic and learn nothing. So the reinforcement is gated by an observed
usefulness $U_b$, measured after the branch has been searched rather than
predicted before it.

Two further commitments separate the model from ordinary tree search.

The unit of traffic is a complete branch. A branch $b$ is a visited path from
the root to a frontier or terminal node, with edge set $P_b$. A deep
observation is evidence about every prefix decision that made the observation
reachable, so one observation credits every edge on its route in a single
event.

The policy head sets initial conditions only. An ordinary move policy
$\pi_\theta(a \mid s)$ supplies a branch mass
\begin{equation}
  m_\theta(b) = \prod_{k=1}^{d} \pi_\theta(a_k \mid s_k),
\end{equation}
which initializes conductivity on newly represented edges, above a positive
exploration floor $\varepsilon$. The policy does not select which computation
happens next, and it does not fix the final topology.

The discrete update is
\begin{equation}
  D_e^{t+1} = \rho\, D_e^{t}
  + \alpha \sum_{b \,:\, e \in P_b} |q_b^{t}| \, U_b^{t},
  \qquad 0 < \rho < 1 .
  \label{eq:update}
\end{equation}
Sections~\ref{sec:init} through~\ref{sec:budget} state what is proved about
this update.

\section{The formal model}

The development works over arbitrary finite types \lean{Edge} and \lean{Branch}
with rational arithmetic. Exact rationals are used throughout, so the
executable diagnostic in Section~\ref{sec:budget} has no floating-point slack.

A branch is represented by its edge set, \lean{path : Branch → Finset Edge}.
The model fixes the bookkeeping of flow, credit and decay without committing to
how $q_b$ and $U_b$ are computed. The four definitions that carry the model are
these.

\begin{Verbatim}[fontsize=\small,xleftmargin=1em]
def branchPrior (localProbabilities : List ℚ) : ℚ :=
  localProbabilities.prod

noncomputable def policyConductivity
    (path : Branch → Finset Edge) (branchMass : Branch → ℚ)
    (explorationFloor : ℚ) (edge : Edge) : ℚ :=
  explorationFloor +
    ∑ branch, if edge ∈ path branch then branchMass branch else 0

noncomputable def depositedTraffic
    (path : Branch → Finset Edge) (flow usefulness : Branch → ℚ)
    (edge : Edge) : ℚ :=
  ∑ branch, if edge ∈ path branch then flow branch * usefulness branch else 0

noncomputable def evolve
    (path : Branch → Finset Edge) (flow usefulness : Branch → ℚ)
    (retention learningRate : ℚ) (conductivity : Edge → ℚ)
    (edge : Edge) : ℚ :=
  retention * conductivity edge +
    learningRate * depositedTraffic path flow usefulness edge
\end{Verbatim}

\lean{evolve} is equation~\eqref{eq:update}, with \lean{retention} for $\rho$
and \lean{learningRate} for $\alpha$.

\section{Policy initialization}
\label{sec:init}

\lean{branchPrior\_nil} gives the empty branch prior $1$.

\lean{branchPrior\_append} gives $m(u \mathbin{+\!\!+} v) = m(u)\, m(v)$. Prefix
masses compose, so an edge created during search can be initialized from its
prefix without enumerating the tree.

\lean{branchPrior\_extension\_le}: if $m(u) \ge 0$ and the appended probability
lies in $[0,1]$, then $m(u \mathbin{+\!\!+} [p]) \le m(u)$. Branch priors are
non-increasing in depth. This is the reason a product prior on its own cannot
carry a deep line, it decays geometrically, and the exploration floor rather
than the prior is what keeps deep edges alive.

\lean{policyConductivity\_pos}: for a strictly positive floor $\varepsilon > 0$
and nonnegative branch masses, every initialized edge has $D_e^0 > 0$.

That positivity result does real work. If flow is proportional to conductivity,
an edge initialized at exactly zero carries no flow, so it receives no deposit,
so it stays at zero. A policy assigning probability zero to a move would then
delete that move from the search permanently. The floor is what keeps the
policy a prior instead of a constraint.

\section{Flow bookkeeping}

\lean{branchFlow\_conservation\_at\_split} takes \lean{active : Finset Branch},
the branches through an internal node, a function
\lean{nextChild : Branch → Child} giving the child each one continues through,
and \lean{flow : Branch → ℚ}. It states
\begin{equation}
  \sum_{c} \Bigg( \sum_{b \in \mathrm{active}}
  [\, \mathrm{nextChild}(b) = c \,] \; q_b \Bigg)
  = \sum_{b \in \mathrm{active}} q_b .
\end{equation}

The hypothesis is worth stating plainly: each active branch continues through
exactly one child, which is what \lean{nextChild} being a function encodes.
The theorem is then the re-indexing identity that outgoing child traffic sums
to incoming traffic. It is conservation in the bookkeeping sense. It is not a
derivation of Kirchhoff's law from a graph Laplacian, and it does not show that
any particular flow solver produces a conserved field.

\section{Whole-branch credit, decay, and amplification}

The single-event form isolates one transported branch.

\begin{Verbatim}[fontsize=\small,xleftmargin=1em]
def reinforceBranch (path : Finset Edge)
    (retention learningRate flow usefulness : ℚ)
    (conductivity : Edge → ℚ) (edge : Edge) : ℚ :=
  retention * conductivity edge +
    if edge ∈ path then learningRate * flow * usefulness else 0
\end{Verbatim}

\lean{reinforceBranch\_on\_path}: every edge on the branch gets the same direct
credit, whatever its depth,
\begin{equation}
  D'_e - \rho D_e = \alpha \, q \, u
  \qquad \text{for all } e \in P_b .
\end{equation}

\lean{reinforceBranch\_off\_path}: every edge off the branch gets decay and
nothing else, $D'_e = \rho D_e$.

\lean{whole\_branch\_credit}, summing over the branch,
\begin{equation}
  \sum_{e \in P_b} \big( D'_e - \rho D_e \big)
  = |P_b| \cdot \alpha \, q \, u .
  \label{eq:credit}
\end{equation}

\lean{total\_depositedTraffic} is the many-branch version, by finite double
counting,
\begin{equation}
  \sum_{e} H_e = \sum_{b} |P_b| \, q_b u_b ,
  \qquad H_e = \sum_{b \,:\, e \in P_b} q_b u_b .
\end{equation}

The factor $|P_b|$ in~\eqref{eq:credit} is easy to misread. A longer branch
absorbs proportionally more total deposit, and that is an accounting
consequence of crediting every edge equally, not evidence that long branches
carry more information. If the length preference is unwanted, replace $U_b$ by
$U_b / |P_b|$, or divide by measured branch cost. What the theorem is good for
is forcing that choice into the open instead of leaving it inside a heuristic.

\lean{depositedTraffic\_nonnegative}: nonnegative flow and nonnegative
usefulness give $H_e \ge 0$.

\lean{evolve\_positive}: with $\rho > 0$, $\alpha \ge 0$, all conductivities
positive, and nonnegative flow and usefulness, \lean{evolve} keeps every
conductivity strictly positive. The update never drives an edge to zero, so
decay suppresses a route instead of deleting it, and a faded route stays
revivable.

\lean{unused\_edge\_decays}: for $D_e > 0$ and $\rho < 1$, an edge receiving no
deposit satisfies $\rho D_e < D_e$ strictly.

\lean{larger\_gain\_amplifies\_relative\_conductivity}: writing one step in
multiplicative form $D'_e = (\rho + g_e) D_e$, if $D_A, D_B > 0$ and
$g_B < g_A$ then
\begin{equation}
  D'_B \, D_A < D'_A \, D_B .
\end{equation}
This is the positive-feedback law. It is stated by cross-multiplication, so no
division and no nonzero-denominator side condition appears, and the common
retention term drops out. The edge with the larger useful-flow gain gains
relative conductivity. It is ordered-field algebra and assumes nothing about
where $g_A$ and $g_B$ came from, so it does not say a larger gain corresponds
to a better move.

\section{Budget as integrated traffic, and the exact diagnostic}
\label{sec:budget}

There is no learned action that spends a unit of budget. Computation is
consumed by transport and evaluation, so the budget is an integral,
\begin{equation}
  B_n = \sum_{t<n} \sum_b c_b^{t} \, |q_b^{t}| .
\end{equation}

\lean{usedBudget\_append\_round}: appending a round adds its cost.

\lean{usedBudget\_monotone}: with nonnegative round costs, accumulated budget is
non-decreasing. Stopping at a measured cost is therefore well defined, and the
procedure is anytime.

\subsection*{The two-branch diagnostic}

\lean{BranchFlowMain.lean} runs the dynamics on the smallest configuration that
can show the intended behaviour. Branch A is two unit-length edges deep, branch
B is one edge, and the policy prefers B. Retention is $\rho = 3/4$, usefulness
is $U_A = 4$ against $U_B = 1/2$, and all arithmetic is over \lean{ℚ}.

\begin{Verbatim}[fontsize=\scriptsize,xleftmargin=0pt]
round=0, D={ aRoot := 2/5,    aDeep := 2/5,    bRoot := 3/5   }, flow(A)=1/5,    flow(B)=3/5,   cost=1
round=1, D={ aRoot := 11/10,  aDeep := 11/10,  bRoot := 3/4   }, flow(A)=11/20,  flow(B)=3/4,   cost=37/20
round=2, D={ aRoot := 121/40, aDeep := 121/40, bRoot := 15/16 }, flow(A)=121/80, flow(B)=15/16, cost=317/80
\end{Verbatim}

At the start $1/5 < 3/5$, so the policy's preferred shallow branch carries three
times the traffic. The first useful observation on A credits both of its edges
equally and raises each to $11/10$. By round two the ordering has reversed,
\begin{equation}
  \frac{121}{80} > \frac{15}{16},
\end{equation}
and the deeper branch that the policy disfavoured carries more flow. Four
theorems pin these exact rationals:

\begin{Verbatim}[fontsize=\small,xleftmargin=1em]
initial_policy_prefers_b
first_useful_observation_credits_both_a_edges
useful_deep_branch_eventually_carries_more_flow
initial_round_cost
\end{Verbatim}

This shows that whole-branch reinforcement can overturn a policy
initialization, on one configuration, with usefulness values chosen by hand
rather than measured from a game tree. It says the mechanism does what it is
described as doing. It says nothing about search quality.

These four theorems use \lean{native\_decide}. See Section~\ref{sec:tcb}.

\section{What is not proved}
\label{sec:notproved}

The gap between the theorems above and a working engine is the whole
engineering problem, so it is listed directly.

\begin{itemize}[leftmargin=1.4em]
  \item Nothing about playing strength. No theorem compares branch-flow search
  with alpha-beta, PUCT MCTS, best-first policy search, or random frontier
  expansion. Whether the coupled flow dynamics beat any of them is empirical
  and is not settled here.

  \item Nothing about the usefulness function. $U_b$ is an arbitrary
  nonnegative rational in every theorem. That a particular $U$ correlates with
  playing strength, or resists manipulation, is unproved, and it is the main
  place where the design can fail.

  \item Nothing about convergence. There is no theorem that the conductivity
  field converges, that it concentrates on principal variations, or that the
  induced root decision approaches the minimax decision.

  \item No flow solver. \lean{branchFlow\_conservation\_at\_split} is a
  re-indexing identity over a given flow assignment. No graph-Laplacian solve
  is formalized, and no result bounds the cost of computing the flow field,
  which the loop needs every round and which comes out of the same budget.

  \item No adversarial backup. Minimax, alpha-beta bounds, and the separation
  between the flow potential $p_v$ and the game value $V(v)$ are part of the
  model but are not formalized here.

  \item Nothing about the length preference. The $|P_b|$ factor
  in~\eqref{eq:credit} is proved to exist. Whether it helps or hurts is not
  addressed.

  \item No claim of novelty over ant-colony or MCTS variants. If the coupled
  flow solve is approximated by sampling one path per round and incrementing
  edge counts, the algorithm degenerates toward existing methods. The theorems
  here hold in both cases, so they do not distinguish them.
\end{itemize}

\section{Trusted computing base}
\label{sec:tcb}

The fifteen theorems in \lean{TrickyProof/BranchFlow.lean} depend on Lean's
three standard axioms and nothing else:

\begin{Verbatim}[fontsize=\small,xleftmargin=1em]
propext, Classical.choice, Quot.sound
\end{Verbatim}

These are the axioms of classical higher-order logic as used throughout
Mathlib. Trusting these results means trusting the Lean kernel, those axioms,
and the statements being the intended ones.

The four theorems in \lean{BranchFlowMain.lean} each carry one extra axiom, of
the form

\begin{Verbatim}[fontsize=\small,xleftmargin=1em]
BranchFlowDemo.<name>._native.native_decide.ax_1_1
\end{Verbatim}

\lean{native\_decide} evaluates a decidable proposition with compiled code and
asserts the result, so it trusts the Lean compiler and runtime on top of the
kernel. That is a strictly larger trusted base. The claims affected are only
the exact rationals in the Section~\ref{sec:budget} table. Replacing
\lean{native\_decide} with kernel-level \lean{decide} or \lean{norm\_num} would
remove the dependency, and is worth doing.

The audit reproduces in full:

\begin{Verbatim}[fontsize=\small,xleftmargin=1em]
lake build
lake env lean AxiomAudit.lean
\end{Verbatim}

There are no occurrences of \lean{sorry} or \lean{admit} in the repository.

\section{Reproducing}

\begin{Verbatim}[fontsize=\small,xleftmargin=1em]
lake exe cache get     # optional: prebuilt Mathlib oleans
lake build             # library and diagnostic
lake exe branch_flow_demo
lake env lean AxiomAudit.lean
\end{Verbatim}

Toolchain is \lean{leanprover/lean4:v4.33.1} with Mathlib \lean{v4.33.1}.

\section{Summary}

The claim supported here is narrow and should be read narrowly. Given any
nonnegative flow assignment and any nonnegative usefulness signal, whole-branch
reinforcement with common decay is coherent: initialization stays positive
above a floor, credit reaches every depth of a transported branch in one event,
unused routes strictly decay without being deleted, larger useful gains
strictly amplify relative conductivity, accumulated cost is monotone, and on
one exact two-branch configuration the dynamics reverse a policy initialization
in two rounds.

Coherence is a precondition. It is not a result about search quality, and
nothing above should be read as one. Whether the loop
\begin{equation}
  \text{policy prior} \rightarrow \text{branch flow} \rightarrow
  \text{information} \rightarrow \text{conductivity} \rightarrow
  \text{new branch flow}
\end{equation}
searches a game tree well is still open. The decisive test is likely the
ablation that removes branch-wide credit in favour of endpoint-only credit. If
crediting the whole transported branch does not improve regret per unit of
cost, the main motivation for the design is absent.

\begin{thebibliography}{9}

\bibitem{nakagaki2000}
T. Nakagaki, H. Yamada, \'A. T\'oth.
\emph{Maze-solving by an amoeboid organism.}
Nature \textbf{407}, 470 (2000).

\bibitem{tero2007}
A. Tero, R. Kobayashi, T. Nakagaki.
\emph{A mathematical model for adaptive transport network in path finding by
true slime mold.}
Journal of Theoretical Biology \textbf{244}(4), 553--564 (2007).

\bibitem{tero2010}
A. Tero, S. Takagi, T. Saigusa, K. Ito, D. P. Bebber, M. D. Fricker,
K. Yumiki, R. Kobayashi, T. Nakagaki.
\emph{Rules for biologically inspired adaptive network design.}
Science \textbf{327}(5964), 439--442 (2010).

\end{thebibliography}

\end{document}
